\documentclass[11pt]{article}

\usepackage[margin=1.6cm]{geometry}
\usepackage{amsmath}
\usepackage{amssymb}
\usepackage{enumitem}
\usepackage{titlesec}
\usepackage{xcolor}
\usepackage{fancyhdr}
\usepackage{array}

\titleformat{\section}{\large\bfseries\color{blue!50!black}}{\thesection}{0.6em}{}
\titlespacing{\section}{0pt}{1.1em}{0.4em}

\setlength{\headheight}{14pt}

\pagestyle{fancy}
\fancyhf{}
\fancyhead[L]{Dice Experiment: Rolling a 7}
\fancyfoot[C]{\thepage}
\renewcommand{\headrulewidth}{0.4pt}

\fancypagestyle{firstpage}{%
  \fancyhf{}
  \fancyhead[L]{Dice Experiment: Rolling a 7}
  \fancyhead[R]{Name: \makebox[4cm]{\hrulefill}}
  \fancyfoot[C]{\thepage}
  \renewcommand{\headrulewidth}{0.4pt}
}

\newcommand{\ans}{\hrulefill}
\newcommand{\anshere}{\par\vspace{0.4em}\noindent\ans\par\vspace{0.4em}}
\newcommand{\answork}{%
  \par\vspace{0.3em}
  \noindent $\Omega = $\hrulefill\par\vspace{0.4em}
  \noindent $E = $\hrulefill\par\vspace{0.4em}
  \noindent $P(E) = $\hrulefill\par\vspace{0.3em}
}
\setlist[enumerate,1]{leftmargin=1.4em, itemsep=0.5em, topsep=0.3em}

\newcolumntype{Z}{>{\centering\arraybackslash}p{1.1cm}}

\begin{document}
\thispagestyle{firstpage}

\begin{center}
  {\Large\bfseries Dice Experiment: Rolling a Total of 7}
\end{center}
\vspace{0.3em}
\noindent In this experiment, you will find out how likely it is to roll a total of $7$ when
you roll two dice --- first by \emph{calculating} it, then by \emph{testing} it yourself.

\section{Step 1: Calculate the Expected Probability}
\noindent When two dice are rolled and their scores are added together, there are $36$ equally
likely outcomes. Complete the table below to find every possible total.

\vspace{0.3em}
\begin{center}
\begin{tabular}{|Z|Z|Z|Z|Z|Z|Z|}
\hline
$+$ & $1$ & $2$ & $3$ & $4$ & $5$ & $6$ \\\hline
$1$ & \rule{0pt}{0.8cm} & & & & & \\\hline
$2$ & \rule{0pt}{0.8cm} & & & & & \\\hline
$3$ & \rule{0pt}{0.8cm} & & & & & \\\hline
$4$ & \rule{0pt}{0.8cm} & & & & & \\\hline
$5$ & \rule{0pt}{0.8cm} & & & & & \\\hline
$6$ & \rule{0pt}{0.8cm} & & & & & \\\hline
\end{tabular}
\end{center}

\vspace{0.4em}
\noindent Now use your table to calculate the \textbf{expected (theoretical) probability} of
rolling a total of $7$.
\answork

\section{Step 2: Carry Out the Experiment}
\noindent Use the tally chart below to record your results as you go: one tally for
\textbf{every} roll, and a second tally only when the total is $7$.

\vspace{0.3em}
\begin{center}
\begin{tabular}{|p{4.3cm}|p{8.8cm}|p{1.6cm}|}
\hline
\textbf{} & \textbf{Tally} & \textbf{Total} \\\hline
Number of $7$'s rolled & \rule{0pt}{1.6cm} & \\\hline
Number of rolls & \rule{0pt}{1.6cm} & \\\hline
\end{tabular}
\end{center}

\vspace{0.4em}
\noindent\textbf{Instructions:}
\begin{enumerate}[label=\arabic*.]
  \item Roll two dice together and add the two scores to find the total.
  \item Add one tally mark to \textbf{``Number of rolls''}.
  \item If the total was $7$, also add one tally mark to \textbf{``Number of $7$'s rolled''}.
  \item Repeat steps 1--3 until you have rolled \textbf{at least $30$ times}.
  \item Count up your tally marks and write each total in the right-hand column above.
\end{enumerate}

\section{Step 3: Calculate Your Experimental Probability}
\[
  P(7)_{\text{experimental}} = \frac{\text{number of } 7\text{'s rolled}}{\text{number of
  rolls}} = \ans
\]
\noindent How does your experimental probability compare to the expected probability you
calculated in Step 1? Why might they be different?
\anshere
\anshere

\section{Step 4: Explanatory and Response Variables}
\noindent In this experiment, which is the \textbf{explanatory variable} and which is the
\textbf{response variable}: the \emph{number of rolls}, or the \emph{number of $7$'s rolled}?
Explain your answer.
\anshere
\anshere

\end{document}
